\documentclass[language=chinese]{campusbeamer}
\usetikzlibrary{arrows.meta,positioning}
\title{流程图布局}
\author{Template check}
\begin{document}
\section{流程图}
\subsection{节点与连接}
\begin{frame}{服务器与本地代理}
\begin{figure}
\centering
\begin{tikzpicture}[
  node distance=1.2cm and 1.7cm,
  flow node/.style={draw=maincolor, fill=maincolor!8, rounded corners=2pt,
    minimum width=3.3cm, minimum height=0.9cm, align=center,
    font=\small\sffamily, text=black},
  flow edge/.style={-{Latex[length=2.2mm]}, draw=maincolor, line width=.7pt}
]
\node[flow node] (local) {本机\\代理};
\node[flow node, right=of local] (ssh) {SSH\\通道};
\node[flow node, right=of ssh] (server) {远程\\服务器};
\draw[flow edge] (local) -- node[above,font=\scriptsize]{请求} (ssh);
\draw[flow edge] (ssh) -- node[above,font=\scriptsize]{转发} (server);
\end{tikzpicture}
\caption{本机代理经 SSH 通道连接远程服务器的示意路径}
\end{figure}
\end{frame}
\end{document}
